% ECONOMICS_PROBLEM: {"id": "E31-439", "title": "Inflation dynamics and expectation formation", "jel_code": "E31", "source_id": "OP-0012"}
% !TeX program = xelatex
\documentclass[11pt,a4paper]{article}
\usepackage[margin=23mm]{geometry}
\usepackage{fontspec}
\setmainfont{Latin Modern Roman}
\usepackage{amsmath,amssymb,mathtools}
\usepackage{xcolor,enumitem,booktabs,longtable,array,xurl,hyperref}
\definecolor{muted}{HTML}{53636C}
\hypersetup{colorlinks=true,urlcolor=blue,linkcolor=blue}
\setlength{\parindent}{0pt}
\setlength{\parskip}{5pt}
\newcommand{\R}{\mathbb R}
\newcommand{\E}{\mathbb E}
\newcommand{\N}{\mathbb N}
\newcommand{\ind}{\mathbf 1}
\newcommand{\Var}{\operatorname{Var}}
\newcommand{\Cov}{\operatorname{Cov}}
\newcommand{\supp}{\operatorname{supp}}
\DeclareMathOperator*{\argmin}{arg\,min}
\DeclareMathOperator*{\argmax}{arg\,max}
\newcommand{\entrymeta}[3]{{\sffamily\small\color{muted}\textbf{\texttt{#1}}\quad|\quad #2\par #3\par}}
\newcommand{\Problemref}[1]{\texttt{#1}}
\newcommand{\JELref}[2]{\texttt{#2} (JEL \texttt{#1})}
\begin{document}
\section*{Inflation dynamics and expectation formation}
\textbf{Problem ID:} \texttt{E31-439}\quad \textbf{JEL:} \texttt{E31}\par
% BEGIN ECONOMICS STATEMENT
\label{problem:OP-0012}
\label{atlas:prob:A2}

\noindent\textbf{Original atlas ID:} A2. \textbf{Formulation:} editorial specification of a research family.

\subsection*{Definitions and assumptions}

Define inflation $\pi_t$ using a declared price index and frequency, real activity gap $x_t$ relative to a specified benchmark, and predecision regime $s_t$. Firm $f$ holds a subjective conditional belief $B_{ft}$ over future inflation given its information; its forecast is $b_{ft}=\int\pi_{t+1}\,B_{ft}(\mathrm d\pi_{t+1})$. Let $b_t=\sum_fw_{ft}b_{ft}$, with nonnegative pricing-relevant weights summing to one. As an illustrative restricted model, posit
\[
\pi_t=\beta b_t+\kappa(s_t)x_t+\gamma(s_t)^\top z_t+\eta_t,
\]
where $\beta$ is specified or estimated, $z_t$ contains measured cost disturbances, and $\eta_t$ is an unobserved pricing disturbance. An expectation-updating model indexed by $\psi$ determines $B_{ft}$; survey answers are measurements of these beliefs under an explicit reporting and sampling model.

\subsection*{Formal research question}

For a stated class of price-setting, belief-updating, and survey-measurement models, characterize the parameters $(\beta,\kappa,\gamma,\psi)$ and shock-response functionals identified by joint micro prices, forecasts, aggregate outcomes, and valid instruments. For example, declared instruments $Z_t$ must satisfy $\mathbb E[Z_t\eta_t\mid s_t]=0$ and the rank condition needed to identify the corresponding coefficients. Determine whether models differing in information rigidity, price adjustment, or expectation formation are observationally distinguishable. Estimate conditional inflation responses to prespecified externally identified supply and demand disturbances, allowing coefficients and information acquisition to vary with the pre-shock state. Compare genuine out-of-sample forecasts and account for parameter and model uncertainty.

\subsection*{Interpretation and scope}

The displayed equation is an editorial benchmark, not a structural consequence of every nominal-rigidity model. Its subjective expectation aggregate need not equal a household survey average, a market-implied expectation, or a rational full-information forecast. The gap benchmark, firm weights, survey selection, and measurement noise materially affect identification. Endogenous activity and beliefs invalidate ordinary least-squares causal interpretations without additional assumptions. Large shocks can change both pricing hazards and information acquisition, so a stable linear coefficient is a testable restriction rather than a maintained truth. Matching firm forecasts to their own pricing records is especially informative if observation and attrition processes are included in the model.

\subsection*{Source audit and status}

Friedman (1968), Mankiw--Reis (2002), and Coibion--Gorodnichenko (2015) are verified with their exact primary records. They provide expectations-based monetary reasoning, a sticky-information alternative, and a forecast-revision framework with empirical evidence. None establishes that a single observed expectation series governs all firms' inflation decisions. The source atlas's broad wording mixes prediction, expectation measurement, and causal pricing mechanisms. These are retained as a restricted empirical identification agenda, not an unresolved Phillips-curve theorem. The formal equation, weighting, regime specification, and validation target above are editorial; any new application must justify its instruments and the interpretation of its belief measures.

\subsection*{References}

\begin{enumerate}

\item[S1] Milton Friedman (1968). \emph{The Role of Monetary Policy}. American Economic Review 58(1), 1--17. \url{https://miltonfriedman.hoover.org/objects/57306/the-role-of-monetary-policy}. \textit{Locator:} presidential address; expectations and natural-rate discussion. \textit{Establishes:} The theoretical role of expectations in monetary-policy limits. \textit{Does not establish:} identified contemporary firm-level expectation mechanisms.

\item[S2] N. Gregory Mankiw and Ricardo Reis (2002). \emph{Sticky Information versus Sticky Prices: A Proposal to Replace the New Keynesian Phillips Curve}. Quarterly Journal of Economics 117(4), 1295--1328. \url{https://mankiw.scholars.harvard.edu/publications/sticky-information-versus-sticky-prices-proposal-replace-new-keynesian-phillips-}. \textit{Locator:} abstract; sticky-information model. \textit{Establishes:} An alternative information-rigidity mechanism for inflation dynamics. \textit{Does not establish:} a universal state-invariant inflation forecasting law.

\item[S3] Olivier Coibion and Yuriy Gorodnichenko (2015). \emph{Information Rigidity and the Expectations Formation Process: A Simple Framework and New Facts}. American Economic Review 105(8), 2644--2678. \url{https://pubs.aeaweb.org/doi/10.1257/aer.20110306}. \textit{Locator:} abstract; forecast-error/forecast-revision framework. \textit{Establishes:} Empirical tests and evidence concerning information rigidity. \textit{Does not establish:} that survey averages are automatically the expectations entering pricing equations.

\end{enumerate}

\subsection*{Common conventions: Source mathematical conventions}

Unless an entry states otherwise, agent and firm sets are finite. State and action spaces are measurable, strategies and outcome maps are measurable, probability laws are specified, and expectations exist for the targets under discussion. Infinite-horizon discount factors lie in $(0,1)$ and discounted objectives are finite. Norms, tolerances and complexity input sizes must be specified for computational comparisons. Constants or primitives that appear locally are inputs, not empirical facts asserted by this catalogue.

\textbf{Identification.} A statistical model $m\in\mathcal M$ implies an observable law $P_m$ and target $\theta(m)$. At law $P$, its identified set is
\[
\Theta(P;\mathcal M)=\{\theta(m):m\in\mathcal M,\ P_m=P\}.
\]
Point identification holds when this set is a singleton, or equivalently when $P_{m_1}=P_{m_2}$ implies $\theta(m_1)=\theta(m_2)$ for all compatible models. Sharp bounds mean bounds attained, or approached when the boundary is not attained, by compatible models. Writing this set is a definition, not a solution: the substantive task is to establish informative restrictions and derive or estimate the set. An unrestricted-confounding model may give only trivial bounds.

\textbf{Inference.} Identification, estimability and valid uncertainty quantification are separate questions. Fix $\alpha\in(0,1)$. A confidence procedure $C_n$ over a declared law class $\mathcal P$ has asymptotic uniform coverage at level $1-\alpha$ when
\[
\liminf_{n\to\infty}\inf_{P\in\mathcal P}\Pr_P\{\theta(P)\in C_n\}\geq1-\alpha.
\]
For set-identified targets the coverage event must be defined separately, for example coverage of the full identified set. If an entry uses identification-indexed models instead of uniquely indexed laws, the same coverage convention is applied over those models. Pointwise limits do not establish uniform validity, and asymptotic validity does not guarantee finite-sample precision. Sample sizes, dependence structures and weak-identification sequences are part of each application.

\textbf{Counterfactuals.} $Y_i(a)$ is a potential outcome under a specified intervention, including its timing, implementation and versions. It is not $Y_i$ conditional on voluntarily choosing $a$. A full intervention vector permits interference; shorthand scalar treatment notation does not silently impose no interference. Assignment, selection, attrition, anticipation and measurement must be specified. A claim about an instrument's local effect is not automatically a population policy effect.

\textbf{Economic equilibrium.} A counterfactual model includes best responses, constraints, feasibility, market clearing, state transitions and expectations consistent with the stated information. When several equilibria exist, either a declared selector is an input or the target is set-valued. Comparisons across policy regimes must use the stated selection convention. Reduced-form relationships are not presumed invariant after an equilibrium-changing intervention.

\textbf{Optimization and welfare.} Optimization is over the stated feasible set. If attainment is not established, $\sup$ or $\inf$ and $\varepsilon$-optimal choices replace an asserted maximizer or minimizer. For example, with value $V=\sup_{a\in\mathcal A}W(a)<\infty$, an $\varepsilon$-optimal set is $\{a\in\mathcal A:W(a)\geq V-\varepsilon\}$ for $\varepsilon>0$. Benefits, costs, risk and health or environmental outcomes can be aggregated only using declared units and conversion weights. Interpersonal utility weights, the treatment of fiscal transfers and foreign profits, discounting, and the behavioral welfare criterion are explicit inputs. A comparative-static derivative additionally requires existence and appropriate differentiability of the selected counterfactual.

\textbf{Sources.} Local labels [S1], [S2], and so forth restart in every question. Locators identify the relevant abstract, model, section or result at the level actually checked; a bibliographic or abstract-level verification is not represented as inspection of an entire proof. Working papers are distinguished from later publications and changing versions. Source summaries are paraphrased. All URL checks and editorial formulations refer to the audit date stated above.
% END ECONOMICS STATEMENT
\end{document}
